\chapitre{Reconstruction en mémoire linéaire}
Une trace optimale passe par une colonne de la ligne médiane \(h=\lfloor m/2\rfloor\). Le principe de coupe de Hirschberg~\cite{hirschberg}, initialement formulé pour une plus longue sous-séquence commune, s'applique aux coûts unitaires d'édition. Une passe avant calcule \(F_j=d(A_{1:h},B_{1:j})\). Une passe en sens inverse calcule \(G_{n-j}=d(A_{h+1:m},B_{j+1:n})\). En coupant un chemin optimal à la ligne \(h\), puis en concaténant deux solutions optimales pour cette coupe, on obtient
\[
 d(A,B)=\min_{0\le j\le n}(F_j+G_{n-j}).
\]
Le noyau retient une colonne minimisante et résout récursivement les deux rectangles. Il place la plus courte chaîne en colonnes et réutilise les deux mêmes tableaux dans les appels enfants. Les scores occupent \(8(\min(m,n)+1)\) octets, auxquels s'ajoutent la pile de profondeur logarithmique et le CIGAR de sortie. À chaque niveau, la somme des aires décroît avec les hauteurs des rectangles. Le travail reste en \(O(mn)\).
\chart{21_trace_lineaire.png}{Sur un cœur du M1 Pro, cinq processus par cas comparent le coût de reconstruction et la mémoire résidente des deux méthodes.}{fig:trace}{bench/traceback.py}
La figure~\ref{fig:trace} utilise des préfixes de la même paire synthétique. Sur le carré de 8\,000 bases, le temps médian passe de 145,51 à 226,31\,ms avec la reconstruction par coupe. Le pic RSS médian passe de 246,70 à 2,55\,Mio. Les deux chemins coûtent 203 éditions.
\synthese{À 25\,000 bases par chaîne, Hirschberg restitue un chemin de 531 éditions en 2,253\,s avec un pic de 2,80\,Mio.}{Quand seul le score est demandé, le calcul peut aussi limiter le nombre de diagonales visitées.}
